\section{System Overview}
\label{sec:overview}

Figure~\ref{fig:architecture} shows the framework as five stages with two
feedback loops. A registered TLS scan enters the \emph{deterministic
geometric backbone} (green in the figure): preprocessing and structure
removal, object decomposition with remnant filtering, and explicit modeling.
Its output---a set of clusters with explicit geometric attributes and
provisional labels---feeds the \emph{probabilistic semantic stages} (orange):
multi-view VLM verification with an internal escalation loop, and
confidence-gated integration into the knowledge graph. The second loop is
the \emph{update loop}: a later scan of the same site re-enters the backbone,
is matched against existing node identities, and applies a node-level upsert
rather than a rebuild.

The color separation in the architecture is a claim boundary, not a drawing
convention. Everything in the backbone is seeded and parameter-fixed;
re-running it on the same host reproduces every retained cluster
byte-identically (measured in Sec.~5.2; across hosts we define reproducibility
numerically, within floating-point tolerance). Everything downstream of the
first VLM call is stochastic; there we claim \emph{control}---schema-forced
outputs, multi-view voting, automatic escalation, and confidence
gating---rather than reproducibility.

Each object node carries the three TOSM description layers: a
\emph{symbolic} model (type, name, status), an \emph{explicit} model
(pose, extents, yaw, color), and an \emph{implicit} model (movability,
openability, key-object status). Verified and unverified content flow through
the same graph but are separated at every interface: verification status is
part of the node schema, gating is applied structurally at query time
(Sec.~\ref{sec:query}), and the place layer is derived from verified objects
only (Sec.~\ref{sec:place}).

The framework is fully automated \emph{after registered TLS input}: scanner
operation and vendor multi-setup registration are manual; everything from E57
to queryable graph---including escalation decisions and update
classification---runs without human intervention, and objects the verifier
cannot resolve remain explicitly \texttt{unverified} rather than being
resolved by hand.

\begin{figure}[t]
\centering
\includegraphics[width=0.95\linewidth]{fig1_architecture.png}
\caption{Framework architecture. Green stages form the deterministic
geometric backbone; orange stages are probabilistic semantic inference under
voting, escalation, and gating control. Two loops: automatic escalation
(within verification) and incremental update (across scans).}
\label{fig:architecture}
\end{figure}
